% RESEARCH_DRAFT_SHA256: bf2997d6577373c3b5b573c18a692ec9e338ea0b0e47f0b86a708dcc7d8d9d78
% Layout and source-reference wrappers expanded for the web renderer.
\documentclass[11pt]{article}
\usepackage[margin=25mm]{geometry}
\usepackage{amsthm}
\usepackage{fontspec,unicode-math}
\setmainfont{Latin Modern Roman}
\setmathfont{Latin Modern Math}
\usepackage{amsmath,amssymb,mathtools,xurl,hyperref}
\hypersetup{colorlinks=true,urlcolor=blue,linkcolor=blue}
\setcounter{secnumdepth}{0}
\setlength{\parindent}{0pt}
\setlength{\parskip}{5pt}
\setlength{\emergencystretch}{3em}
\newcommand{\sref}[2]{\hyperlink{source:#1:#2}{[S#2]}}
\newcommand{\cataloguescope}[1]{\paragraph{Source record.}#1}
\newcommand{\floor}[1]{\left\lfloor#1\right\rfloor}
\newtheorem{proposition}{Proposition}
\newtheorem{lemma}[proposition]{Lemma}
\begin{document}
\paragraph{Research manuscript.} OpenAI \texttt{gpt-6-astra}, reasoning \texttt{ultra}; 11 October 2026.
\paragraph{Preserved source body.} The complete catalogue mathematical text and sources follow with layout and source links expanded, including source qualifications and its historical preparation notes. New research begins at the section headed Research attempt.

\section{An event access assurance case for mathematics}\label{statement}

\subsection{Definitions and mathematical statement}

\paragraph{Shared notation.}Standard number systems are $\mathbb N,\mathbb Z,\mathbb Q,\mathbb R,\mathbb C$. Specialized notation and index ranges follow the cited source and the definitions below.


An accessible event allows intended participants to use its venues, materials, schedules, and communication channels. Participation barriers can concern mobility, caregiving, travel, finances, identity-related needs, or connectivity. Organizer resources can document arrangements, responsibilities, and processes for obtaining accommodations.
\paragraph{Record scope and qualifications.}
The catalogue title is retained as metadata. Named certificates, contracts, bounds or frameworks in that title are not asserted as results of the original source extract.
The source registry labels formal reconstruction unrecoverable. The supplied agenda is shown literally, including its unfinished participant list; no formal optimization problem or universal accessibility standard is inserted.

Catalogue screening limitation: The supplied source audit marks the record unrecoverable: No word is inserted after the comma: the missing category is unrecoverable from the sources checked. The verified named groups are parent participants, trans participants, participants from the Global South, and non-US participants. Disability and universal design, caregiving beyond parenting, language, time zones, religion and culture, cost, digital access, health, and safety are treated below as a present-day expansion needed to answer the clear general question, not as reconstruction of the missing text.

The source snapshot labels this record partially solved; this is catalogue metadata rather than independent verification of a mathematical result or the current completion of a project.
The source text below is rendered literally where necessary; damaged formulas and composite extracts are not silently repaired into a new mathematical claim.
The supplied catalogue marks this record as unrecoverable. Its extract and terminology are documented, but a complete original problem has not been reconstructed.
\paragraph{Statement recorded in the supplied source.}
Improving Conference Organization/Resources for Organizers

What might it look like to make our events (conferences, workshops, etc.) accessible to all participants? How can we make it accessible to parent participants, trans participants, participants from the Global South, non-US participants, ? Can we make a "best practices" document for anyone interested in organizing an equitable event in math? \href{http://aimpl.org/gemscombin/7/}{[S2]}
\paragraph{Source status.}
The supplied catalogue labels this source record partially solved; that metadata was checked against the supplied snapshot on 2026-10-11. The label is preserved for provenance and is not an independent verification of present open status or mathematical correctness.
\subsection{Short English statement}
Preserve the proposal for more accessible mathematical events and resources for organizers.
\subsection{Sources}
\begin{itemize}
\item[S1] Ulam AI and the UnsolvedMath Contributors, UnsolvedMath: A Curated Collection of Open Mathematics Problems, catalogue entry. Catalogue curation and metadata: CC BY 4.0. The source extract is retained; definitions, English summaries and qualifications are editorial additions.
\url{https://huggingface.co/datasets/ulamai/UnsolvedMath}.
\item[S2] AIM Workshop Problem Lists, source linked by the catalogue.
\url{http://aimpl.org/gemscombin/7/}.
\end{itemize}
\paragraph{Preparation and review.}
The original source and UnsolvedMath contributors retain attribution for the source extract. Codex (GPT-6; exact serving revision not exposed) assisted with definitions, English summaries, scope notes and formatting. No new human mathematical authorship or independently verified solution is asserted. Catalogue identity checked against the complete supplied registry and current repository. Distinct upstream records, including repeated questions and broad or damaged extracts, are retained for complete coverage of the source CS category. A repeated record is not a new independent conjecture. Complete mathematical-equivalence and current-literature review has not been performed. Preparation and final review follow contribution prompts C and D, with source gaps disclosed. Structural and rendering checks do not establish a new theorem or open status.
\label{end:statement}

\section{Research attempt}
\subsection{Literature, source context and scope}
The complete catalogue body above retains the original agenda, source status, and historical preparation. The AIMPL problem page could not be retrieved. The supplied complete source transcription was read, and the accessible primary AIM workshop announcement was checked for context. The announcement is not a substitute for verification of the inaccessible problem page. The catalogue's earlier proposed frameworks and research summaries are context, not established solutions credited to this attempt. The following is an explicitly limited investigation of one route; it neither substitutes a formalized toy model for the entire source agenda nor claims empirical or deployed-software results. Elementary auxiliary facts are not new original-target progress.

\subsection{Testing a feature-count measure of event accessibility}
The source asks about many participants and many barriers. We isolate a limited journey model: a participant can attend if a directed graph of usable steps contains a path from start $s$ to successful participation $t$. A step is usable or blocked for this modeled participant. These binary labels are stipulated, not inferred from a person's identity, and do not represent all dimensions of inclusion.
\begin{proposition}
A fraction of usable features arbitrarily close to one does not imply an accessible end-to-end journey. Removing a set of edges $B$ prevents attendance exactly when every original $s$--$t$ path contains an edge of $B$.
\end{proposition}
\begin{proof}
For the first assertion take a chain of $m\ge2$ steps and block exactly one. The usable fraction is $(m-1)/m$, tending to one as $m$ grows, but the unique start-to-finish path is broken. For the second, if an original path avoids $B$, all its edges remain and it still permits attendance. Conversely every path in the remaining graph is also an original path avoiding $B$; if every original path meets $B$, no such remaining path exists.
\end{proof}
An explicit finite test explores usable out-neighbors from $s$ until no new vertex is found. Every visited vertex has a usable path by induction on discovery, and every usable path endpoint is eventually visited by induction on path length. Thus $t$ is visited exactly when a usable journey exists. This checks paths for specified personas and conditions; it cannot certify accessibility for all real participants unless those conditions are complete and accurate.
\subsection{Verification and remaining gaps}
A 100-edge chain with edge 50 blocked has 99 usable steps and no complete path, as a traversal verified. The attempted inference from a high checklist percentage to universal event access fails. No actual event audit, participant consultation, organizer resource document, or coverage of the source's unfinished list of participant needs was completed. The original broad agenda remains intact.

\section{References and source attribution}
Primary source: \url{http://aimpl.org/gemscombin/7/}. Attempted primary AIMPL page retrieval; failed with timeout or gateway error. The supplied full source transcription was read; the inaccessible live page was not independently inspected.

\url{https://aimath.org/pastworkshops/gemscombin.html}. Read the primary March 2023 workshop announcement and its community inclusion agenda.

All retrievals were checked on 11 October 2026. Original contributor and catalogue attribution remains in the preserved Sources section. The elementary arguments below are self-contained; their exposition is not a claim of first discovery.


\paragraph{Integration review.} Codex (GPT-6; exact serving revision not exposed) reviewed the written arguments and expanded layout and source-reference wrappers. Repository metadata and the production web renderer were checked using Python 3.13 and JavaScript with MathJax 3.2.2. This model-assisted review is not independent human proof verification; no PDF compilation or formal proof checking is asserted.

\section{Completion Estimate}
\noindent\textbf{COMPLETION: 0\%}
No substantiated new progress toward the entire original target is claimed. The worked results and failed route are delimited above; a known result, illustrative model, or auxiliary obstruction is not counted as original-target completion.

\section{Final status and attribution}
\textbf{Attempt status: unresolved.} This describes this research attempt, not a new ruling on the source catalogue's status or on external literature claims.
The mathematics and this self-audit were generated by OpenAI \texttt{gpt-6-astra}, reasoning \texttt{ultra}, on 11 October 2026. Codex (GPT-6; exact serving revision not exposed) orchestrated the batch and contributed formatting/repository review. No human mathematical authorship, independent human proof checking, or formal verification is claimed. The historical preparation paragraphs in the reproduced catalogue body concern that body only.

Contribution prompt A was applied to the complete definition and supplied source record, followed by the adversarial checks of prompt B. Local reproduction files are \texttt{batch013/checks.py}, \texttt{computations.json}, and the per-ID audit JSON. They were not public links at preparation time. Numerical checks are finite and do not prove unrestricted results. No LaTeX engine was found locally; compilation was not performed by this research agent.

\end{document}
